\documentclass{cv}

\usepackage[utf8]{inputenc}
\usepackage[ngerman]{babel}
\usepackage{hyperref}

\hypersetup{
	colorlinks=true,
	urlcolor=myblue,
	linkcolor=myblue,
	citecolor=myblue,
	pdftitle={Dr. André Kaderli -- Lebenslauf (Dozentur Quantencomputing)},
	pdfauthor={André Kaderli}
}

% Rote Platzhalter \tbd{...} vor dem Versand ausfuellen oder entfernen.

\begin{document}

	\name{Dr. André Kaderli}
	\tagline{Promovierter theoretischer Physiker (summa cum laude) mit Ausbildung in Quantentheorie und mathematischer Physik an der ETH Zürich und über fünf Jahren Praxis in Machine Learning und mathematischer Modellierung -- damit bringe ich beide Grundlagen des Quantum Machine Learning mit. Mit über 15 Jahren Lehrerfahrung verbinde ich anwendungsorientierte Forschung und verständlichen Mathematikunterricht.}
	\begin{contactline}
		Obere Bahnhofstrasse 22M, 3700 Spiez\sep
		\href{mailto:andre.kaderli@gmail.com}{andre.kaderli@gmail.com}\sep
		\href{tel:+41793323931}{+41\,79 332 39 31}\sep
		\href{https://www.linkedin.com/in/andre-kaderli/}{linkedin.com/in/andre-kaderli}
	\end{contactline}

	\datedsection{Profil für Quantencomputing und QML}{}
	\kv{Quantentheorie:}{Quantenmechanik, Quantenfeldtheorie, Statistische Physik; Quanteninformation und Quantencomputing (Vorlesung \tbd{Kursname}, ETH Zürich)}
	\kv{Mathematische Grundlagen:}{Lineare Algebra und Hilberträume, Gruppen- und Darstellungstheorie, Wahrscheinlichkeitstheorie, Differentialgeometrie, Numerik und Optimierung}
	\kv{Machine Learning in der Praxis:}{Prognose- und Klassifikationsmodelle im produktiven Einsatz, statistisches Lernen, Anomalieerkennung, Python-Ökosystem}
	\kv{Forschungsinteressen:}{Variationale Quantenalgorithmen und Quantum Kernels; Vergleich mit klassischen ML-Verfahren an realen Datensätzen aus Verwaltung und Wirtschaft; Quantenalgorithmen für Optimierung und Simulation}

	\datedsection{Lehrgebiete}{}
	\kv{Mathematik:}{Analysis, Lineare Algebra, Wahrscheinlichkeitsrechnung und Statistik, Numerik, Optimierung, Modellierung}
	\kv{Quanten- und Datenwissenschaften:}{Quantencomputing und QML, Machine Learning, Datenanalyse}

	\datedsection{Lehre und Betreuung}{}

	\datedsubsection{Data Academy, Novalytica}{Seit Mai 2022}
	\desc{Kurse zu Statistik, Programmierung (Python, R) und Datenanalyse für Fachleute; rund \tbd{X} Kurse pro Jahr}

	\datedsubsection{Vorlesungsunterstützung, HU Berlin}{2018 -- 2021}
	\desc{Funktionentheorie und Quantenfeldtheorie (\tbd{X} Semester); Betreuung wissenschaftlicher Arbeiten}

	\datedsubsection{Teaching Assistant, ETH Zürich}{2014 -- 2016}
	\desc{Übungsgruppen in Linearer Algebra am Departement Mathematik (\tbd{X} Semester)}

	\datedsubsection{Privater Nachhilfelehrer}{2009 -- 2018}
	\desc{Mathematik, Physik und Informatik}

	\datedsection{Forschung}{}

	\datedsubsection{Max-Planck-Institut für Gravitationsphysik, Potsdam}{}
	\datedsubsubsection{Wissenschaftlicher Mitarbeiter (Doktorand)}{Sep. 2018 -- Aug. 2021}
	\begin{itemize}
		\item Quantengravitation an der Schnittstelle von Stringtheorie und Zahlentheorie (elliptische Multi-Polylogarithmen)
		\item Analytische und numerische Methoden, rekursive Algorithmen zur Lösung komplexer Differentialgleichungen
		\item Konferenzbeiträge und Vorträge: \tbd{Anzahl, ggf. wichtigste Konferenzen}
	\end{itemize}

	\datedsubsection{Ausgewählte Publikationen}{}
	\desc{\tbd{Autoren, «Titel», Zeitschrift Band (Jahr), Seiten}}
	\desc{\tbd{Autoren, «Titel», Zeitschrift Band (Jahr), Seiten}}
	\desc{Vollständige Liste: \href{https://inspirehep.net/authors/1809024}{inspirehep.net/authors/1809024}}

	\newpage
	\datedsection{Ausbildung}{}

	\datedsubsection{PhD in Physik, Humboldt-Universität zu Berlin}{Sep. 2018 -- Aug. 2021}
	\desc{Theoretische Physik (summa cum laude); Dissertation: «\tbd{Titel}», Betreuung: Prof.~\tbd{Name}}

	\datedsubsection{MSc in Physik, ETH Zürich}{Sep. 2016 -- Aug. 2018}
	\desc{Theoretische und Mathematische Physik (Abschlussnote: 6.0/6.0); Fokus auf Quantenfeldtheorie, Allgemeine Relativitätstheorie und Statistische Physik}

	\datedsubsection{BSc in Physik, ETH Zürich}{Sep. 2013 -- Aug. 2016}
	\desc{Physik, Mathematik und Informatik (Abschlussnote: 5.7/6.0)}

	\datedsection{Berufspraxis}{}

	\datedsubsection{Novalytica AG, Bern}{}

	\datedsubsubsection{Principal Data \& AI Engineer und Partner}{Seit Jan. 2026}
	\begin{itemize}
		\item Fachliche Leitung von Daten- und KI-Initiativen; Führung interdisziplinärer Teams (\tbd{X} Personen)
		\item Beratung von Entscheidungsträgern; Übersetzung von Praxisfragen in quantitative Lösungsansätze
	\end{itemize}

	\datedsubsubsection{Senior Data Scientist}{Jan. 2023 -- Dez. 2025}
	\begin{itemize}
		\item Mathematische Modellierung, statistische Analysen und Leitung von ML-Projekten von der Problemformulierung bis zum produktiven Betrieb
	\end{itemize}

	\datedsubsubsection{Data Scientist}{Mai 2022 -- Dez. 2022}
	\begin{itemize}
		\item Entwicklung von ML-Modellen und statistischen Auswertungen; Data-Analytics-Projekte
	\end{itemize}

	\datedsubsection{Eraneos AG, Bern}{}
	\datedsubsubsection{Consultant}{Sep. 2021 -- Apr. 2022}
	\begin{itemize}
		\item IT- und Datenprojekte im öffentlichen Sektor mit Fokus auf quantitative Datenauswertungen
	\end{itemize}

	\datedsection{Weitere Qualifikationen}{}
	\kv{Programmierung:}{Python, R, C++, MATLAB, Mathematica, SQL; \tbd{prüfen: Qiskit/PennyLane}}
	\kv{Zertifizierungen:}{Microsoft Azure, Fabric und Power BI (5\,$\times$ Associate); Databricks Data Engineer Professional; Hermes Advanced}

	\datedsection{Weitere Angaben}{}
	\kv{Sprachen:}{Deutsch (Muttersprache), Englisch (verhandlungssicher), Französisch (Grundkenntnisse)}
	\kv{Stiftungen \& Stipendien:}{Alumni der Schweizerischen Studienstiftung, ehem. Träger des Ernst-Göhner-Stipendiums}
	\kv{Freiwilliges Engagement:}{Tourenleiter beim Schweizer Alpen-Club (SAC)}
	\kv{Interessen:}{Bergsteigen, Klettern, Skifahren, Gleitschirmfliegen, Jiu-Jitsu}

\end{document}
